\section{Conclusion}\label{sec:conclusion}

A particle filter for terrain-aided navigation was implemented, validated against the exact
Kalman filter on planar terrain, and evaluated over real terrain of Sri Lanka in seven Monte Carlo
experiments. The main conclusions are:

\begin{enumerate}
  \item \textbf{The particle filter is the right tool, and real terrain makes the case more
        strongly than synthetic terrain.} Over real hills the EKF already locks falsely in
        \EOneHillSOneZeroZeroEkfFl\,\% of runs at 100\,m initial error and converges in only
        \EOneHillSFiveZeroZeroEkfConv\,\% at 500\,m (synthetic hills:
        \EOneSynthSFiveZeroZeroEkfConv\,\%); this ungated EKF then ends worse than dead reckoning.
        With 20\,000 particles the particle filter still converges in
        \EOneHillSOneZeroZeroZeroPfTwoZeroZeroZeroZeroConv\,\% of the runs at 1\,km and tracks at
        \EOneHillSFiveZeroZeroPfTwoZeroZeroZeroZeroTrack--\EOneKandySFiveZeroZeroPfTwoZeroZeroZeroZeroTrack\,m
        over the hills.
  \item \textbf{Terrain decides what is possible.} Performance follows the terrain information,
        strongly between terrain types and weakly between flights over the same hilly region; over
        flat plains the particle filter stays honestly uncertain instead of locking falsely.
  \item \textbf{False lock is the failure to design against.} More particles help, but roughening
        and a divergence monitor that needs no ground truth are complementary and reach
        \ESixHillSTwoZeroZeroZeroPfMonitorRoughConv\,\% convergence at 2\,km initial uncertainty
        with 5\,000 particles. The resampling scheme and the ESS threshold, by contrast, made no
        significant difference.
  \item \textbf{Robustness and efficiency.} A Student-$t$ likelihood removes most of the
        sensitivity to (zero-mean, 100\,m) altimeter outliers at no cost; constant map or altimeter
        offsets of 10\,m and more remain the most damaging model error; Rao-Blackwellising the INS
        bias fixes the plain filter's impoverished bias estimate once the position is locked on,
        at no extra runtime.
  \item \textbf{How far from optimal.} The unconditional error never undercuts the ensemble
        posterior Cramér--Rao bound; in the tracking phase the 20\,000-particle filter stays about
        twice above the per-flight bound, almost entirely because of its deliberately inflated
        process noise --- a trade made for robust acquisition, where the bound, blind to ambiguity,
        is not informative.
\end{enumerate}

Every number, table and figure in this report is regenerated from saved results by scripts in the
repository, and the study regions and analysis rules were fixed before the experiments ran.
